% ---------------------------------------------------------------------------
\section{Directed microjet: axisymmetric boundary-integral calculation}\label{sec:jet}
This section answers handoff priority 4 (fast mass and focusing) for benchmark B1. A spherical model
cannot select a direction; here the collapse near a rigid receiver is solved in axisymmetric potential
flow with the PFC content laws of Sections~\ref{sec:thermo}--\ref{sec:spherical}.

\subsection{Step J1 --- Governing equations and nondimensionalisation}
\emph{Assumptions.} Inviscid, incompressible, irrotational carrier; rigid plane receiver $z=0$;
uniform cavity pressure $p_b(V)$ (homobaric content); surface tension $\sigma$ on the wall; gravity
negligible ($\rho g\Rm/\Delta p\sim10^{-6}$, with $g$ gravitational acceleration). \emph{Symbols.} $\phi$ velocity potential
($\mathbf u=\nabla\phi$); $S$ the cavity surface; $\mathbf n$ unit normal on $S$ pointing into the liquid;
$q=\partial\phi/\partial n$; $\kappa$ total curvature of $S$ ($\kappa=2/R$ for a sphere);
$D/Dt$ the derivative following a surface node moving with the liquid velocity. Units: length $L$
(the free-field, inviscid, tension-free maximum radius of the same cavity), pressure
$\Delta p_{\mathrm{ref}}=\pinf-\psat(T_0)$, speed $\sqrt{\Delta p_{\mathrm{ref}}/\rho_{\mathrm w}}$, time
$L\sqrt{\rho_{\mathrm w}/\Delta p_{\mathrm{ref}}}$; $\mathrm{We}=\Delta p_{\mathrm{ref}}L/\sigma$.
In the next displayed system, $\mathbf x=(r,z)$, $t$, $\phi$, $\mathbf u=\nabla\phi$, $S$,
$\kappa$ and the spatial derivatives are \emph{dimensionless versions after this scaling};
the dimensional pressures $\pinf,p_b,\Delta p_{\mathrm{ref}}$ (Pa) enter only through their ratio.
$V$ in the content law denotes the physical cavity volume (m$^3$), $\mathrm{We}$ is dimensionless,
$\nabla$ is the gradient, $\nabla^2$ the Laplacian, $|\cdot|$ the Euclidean norm, and $D/Dt$
the material derivative in scaled time. The next equations have consistent scaled units:
\begin{align}
\nabla^2\phi&=0\ \text{in the liquid},\qquad \partial_z\phi=0\ \text{on }z=0,\qquad \phi\to0\ \text{as }|\mathbf x|\to\infty,\label{eq:laplace}\\
\frac{D\mathbf x}{Dt}&=\nabla\phi,\qquad
\frac{D\phi}{Dt}=\tfrac12|\nabla\phi|^2+\underbrace{\frac{\pinf-p_b(V)}{\Delta p_{\mathrm{ref}}}}_{\text{drive}}+\frac{\kappa}{\mathrm{We}}\quad\text{on }S.\label{eq:bern}
\end{align}
\emph{Derivation of \eqref{eq:bern}.} Unsteady Bernoulli in the liquid with the far field at rest:
$\partial_t\phi+\tfrac12|\nabla\phi|^2+p/\rho=\pinf/\rho$. At the wall the liquid pressure is
$p=p_b-\sigma\kappa$ (Laplace). Following a node, $D\phi/Dt=\partial_t\phi+\nabla\phi\cdot\nabla\phi$; substituting
$\partial_t\phi$ gives $D\phi/Dt=\tfrac12|\nabla\phi|^2+(\pinf-p_b+\sigma\kappa)/\rho$; dividing by
$\Delta p_{\mathrm{ref}}/\rho$ gives \eqref{eq:bern}. \emph{Check:} for a constant-pressure cavity the drive
is 1 and \eqref{eq:bern} is exactly the model of Supponen et al.\ (2016, Eq.~5.4, verified).

\subsection{Step J2 --- Boundary integral equation with the receiver image}
\emph{Symbols.} $\mathbf y^*=(y_r,-y_z)$ the mirror image of $\mathbf y$ in the receiver;
$\Gf(\mathbf x,\mathbf y)=|\mathbf x-\mathbf y|^{-1}+|\mathbf x-\mathbf y^*|^{-1}$.
Because $\partial\Gf/\partial z_y=0$ on $z=0$, Green's second identity over the liquid (bounded by $S$, the plane and a
large hemisphere whose contribution vanishes as $\phi=O(|\mathbf x|^{-1})$) involves $S$ only. For $\mathbf x$ on a
smooth part of $S$ the free-term is $2\pi\phi(\mathbf x)$.
For this integral, $\mathbf x,\mathbf y$ are dimensionless field/source coordinates,
$\phi$ dimensionless potential, $q$ dimensionless normal derivative, $\Gf$ the scaled Green kernel,
$S$ the scaled surface, $\dd S_y$ its scaled source area, and
$\partial/\partial n_y$ differentiation along the source unit normal; both integrals use the same
scaled convention as Eq.~\eqref{eq:bern}:
\begin{equation}
2\pi\phi(\mathbf x)+\int_S \Gf(\mathbf x,\mathbf y)\,q(\mathbf y)\,\dd S_y-\int_S\phi(\mathbf y)\,\frac{\partial\Gf}{\partial n_y}(\mathbf x,\mathbf y)\,\dd S_y=0 .\label{eq:bie}
\end{equation}
(Signs: $\mathbf n$ points into the liquid, i.e.\ out of the region bounded by $S$; for a sphere of radius 1 with
$\phi=-1$, \eqref{eq:bie} returns $q=1$, which the implementation reproduces to $1.4\times10^{-8}$ at 64 nodes.) In
axisymmetry the azimuthal integrals reduce to complete elliptic integrals (Appendix~\ref{app:kernels}); the
singular single-layer kernel has a logarithmic singularity integrated with an 8-point log-weighted Gauss rule,
and the free-space double-layer diagonal is fixed by the rigid-mode identity (a uniform $\phi$ produces no
normal flux in the free problem). Interior potentials (for the jet velocity field) follow from the same
representation with $2\pi$ replaced by $4\pi$ and $\mathbf x$ in the liquid.

\subsection{Step J3 --- Discretisation, time integration and stopping rule}
$N+1$ nodes from the far pole ($r=0$, top) to the near pole ($r=0$, bottom); cubic-spline geometry and
potential in chord length with axis symmetry conditions; linear shape functions for $\phi,q$; 10-point
Gauss quadrature on regular elements and 20-point on elements adjacent to the collocation node; classical
RK4 time stepping with $\Delta t=\min(0.02/\max|D\phi/Dt|,\ 0.15\,\Delta s_{\min}/\max|\mathbf u|)$ (Blake's potential
increment control plus a Courant condition); five-point Longuet-Higgins--Cokelet smoothing and
curvature-weighted arc-length redistribution every five steps. A run is declared a jet impact when the
far-pole jet tip comes within $0.01L$ of the opposite cavity wall (the potential-flow solution becomes
singular at contact); the jet speed $U_{\mathrm j}$ is the maximum far-pole speed toward the receiver
before that instant (observable O1). Runs that rebound without impact or produce a non-physical state are
reported as such and not used.

\subsection{Step J4 --- Validation against the constant-pressure benchmark}
Table~\ref{tab:bimval} compares with Supponen et al.\ (2016), Table~1 (rigid-surface column, growth
included, verified from the full text). The BIM jet speeds lie $\bimDevMin$ to $\bimDevMax$\,\% from the
published values for $1.25\le\gamma\le4$; refining from $N=64$ to $96$ changes them by at most
$\bimConvMax$\,\%; halving the potential-increment cap changes them by $\le\bimSensDt$\,\% and halving the initial radius of
the growth proxy by $\le\bimSensRinit$\,\%. The systematic $-4$ to $-6$\% offset is consistent with our
earlier stopping gap ($0.01L$) and our slightly smaller maximum radius in the presence of the wall
($0.99L$ versus their definition with $R_{\max}$). Starting at maximum radius without the growth phase
increases the jet speed by a factor $\bimStartMaxRatio$ at $\gamma=2$, confirming Supponen et al.'s
statement that the growth history strongly affects the jet; the PFC runs therefore include the
adopted real-EOS growth from the assumed activated droplet. This comparison benchmarks the numerical
potential-flow components; it does not experimentally validate a laser-heated PFC microjet.
\input{generated/table_bimval.tex}

\subsection{Step J5 --- PFC-specific content laws}
The PFC cavity grows from the activated droplet at rest ($R_{\mathrm a}/L=0.141$ for PFP) along the
real-fluid isentrope Eq.~\eqref{eq:isentrope}. After the volume maximum four collapse closures
sample plausible limits of the unknown condensation physics; they are not proved bounds on every
aspherical thermodynamic history:
\begin{description}[leftmargin=0.6cm,style=sameline]
\item[frozen] $p_b=p_{\max}(V_{\max}/V)^{k_{\mathrm v}}$, no condensation, adiabatic ($k_{\mathrm v}=\PFPgammav$);
a no-condensation cushion control.
\item[KM $\alpha_{\mathrm c}$] $p_b(V/V_{\max})$ taken from the spherical Keller--Miksis collapse with HKS
condensation coefficient $\alpha_{\mathrm c}=0.1$ or $1$ (Section~\ref{sec:spherical}), continued as
$p_b\propto V^{-1}$ beyond the end of the spherical record (``spherical-content closure'': it assumes the
content pressure--volume history of the aspherical collapse equals that of the spherical one).
\item[isoeq] instantaneous equilibrium at $T_0$: $p_b=\min[\psat(T_0),\psat(T_0)V_{\mathrm{sat}}/V]$; the
equilibrium-at-$T_0$ control, producing a strong drive in these samples.
\end{description}
Surface tension acts throughout ($\mathrm{We}=\PFPWe$ for PFP, $\PFHWe$ for PFH).

\subsection{\texorpdfstring{Results: microjet speeds of a 6\,$\mu$m PFC droplet}{Results: microjet speeds of a 6 micrometre PFC droplet}}
Table~\ref{tab:bimpfc} and Fig.~\ref{fig:bim}. Speed unit $\sqrt{\Delta p_{\mathrm{ref}}/\rho_{\mathrm w}}=\PFPUunit$\,m/s
(PFP), $\PFHUunit$\,m/s (PFH).
\begin{enumerate}[leftmargin=*]
\item \textbf{Largest sampled in-screen jet.} Among the sampled cases inside the Mach-0.3 screen the fastest microjet is
$\PFPenvUmax$\,m/s for PFP (\PFPenvMaxClos, $\gamma=\PFPenvMaxg$) and $\PFHenvUmax$\,m/s for PFH (\PFHenvMaxClos,
$\gamma=\PFHenvMaxg$). Fast-condensation sampled trends straddle the adopted
screen $0.3c_{\mathrm w}=\Ulim$\,m/s. Original-grid log-linear interpolation places it near
$\gamma\approx\PFPcrossIsog$ for PFP isoeq, $\PFHcrossIsog$ and $\PFHcrossKMoneg$ for PFH
isoeq and KM $\alpha_{\mathrm c}=1$; refinement evidence is stated below.
The physical fast-condensation maximum remains unresolved: interpolation is not proof of a
continuous optimum or a physical $U_{\max}\ge\Ulim$ lower bound.
The out-of-screen incompressible values at $\gamma=4$ ($\PFPenvOverMach$\,m/s PFP,
$\PFHenvOverMach$\,m/s PFH, isoeq) are extrapolation estimates, not certified physical upper
bounds. Sound speed is a compressibility scale, not a universal jet-speed cap.
\item \textbf{Condensation is a first-order control.} At fixed $\gamma$ the ordering is
isoeq $\ge$ KM $\alpha_{\mathrm c}=1$ $>$ KM $\alpha_{\mathrm c}=0.1\approx$ frozen in the admissible sampled comparisons.
The largest slow-condensation samples are $\PFPenvSlowUmax$\,m/s (PFP) and
$\PFHenvSlowUmax$\,m/s (PFH); this is not a proof for every $\alpha_{\mathrm c}\le0.1$.
The frozen cushion prevents jet impact
altogether at $\gamma\ge3$: the vapour stops the collapse before the jet can cross the cavity.
\item \textbf{PFP versus PFH.} Per unit $\sqrt{\Delta p_{\mathrm{ref}}/\rho}$, the condensing PFP cavity jets faster than
the constant-pressure benchmark (isoeq at $\gamma=2$: $U^*_{\mathrm j}=20.8$ versus 15.2) because its vapour is
unsaturated at $T_0$ during most of the collapse ($V>V_{\mathrm{sat}}$), so the drive exceeds 1; but its reference
pressure is small ($\PFPdpref$\,kPa), so in physical units PFH (reference $\PFHdpref$\,kPa) gives similar or faster
jets. With the frozen cushion both species give the same 57--65\,m/s at $\gamma\le2$.
\item \textbf{Energy partition.} The jet carries $\PFPenvEfmin$--$\PFPenvEfmax$\% (PFP) and $\PFHenvEfmin$--$\PFHenvEfmax$\%
(PFH) of $\Phia$. At the mass of the fastest admissible jet ($\PFPenvMaxmj$\,pg for PFP) the rigorous energy ceiling
gives a mass-averaged energetic scale of $\PFPenvMaxUceil$\,m/s. The far-pole speed is not a
plug/RMS speed; the rigorous comparison is the profile's kinetic energy against $\Phia$.
\end{enumerate}
\begin{figure}[ht]\centering
\includegraphics[width=0.97\textwidth]{figures/bim_pfc_jets.pdf}
\caption{Microjet speed before impact for a fully activated 6\,$\mu$m droplet versus stand-off, for four
condensation closures. Dotted: the Supponen power law $0.9\zeta^{-1}$ for a constant-pressure cavity scaled with
the same $\Delta p_{\mathrm{ref}}$; dash-dot: Mach 0.3 in water.}\label{fig:bim}
\end{figure}
\input{generated/table_bimpfc.tex}

\subsection{Exact-input PFC refinements and screen-crossing uncertainty}\label{sec:bim_refine}
Nine selected content/stand-off cases were refined from the original control $N=64$ to $N=96$
at the same dimensionless potential-increment cap $\Delta\phi_{\max}=0.02$, and separately
from cap 0.02 to 0.01 at $N=64$. The real-EOS growth history, content arrays, wall tension,
growth seed, smoothing, quadrature and $0.01L$ impact gap are retained; physical-control
comparisons are recorded in \texttt{verification/pfc\_bim\_input\_checks.json}.
Table~\ref{tab:bimpfcconv} retains every requested attempt. Of nine spatial pairs,
$\bimPFCValidSpatialCount$ complete and $\bimPFCAbortedSpatialCount$ PFH $N=96$, 0.02 runs
abort during growth. Their raw pole speeds are diagnostics, not impact predictions.
Among the completed pairs, the largest speed change is $\bimPFCSpatialMaxDeltaU$\%
(PFH KM $\alpha_{\mathrm c}=1$, $\gamma=4$, outside the Mach screen). All nine time-control
pairs complete with changes no larger than $\bimPFCTemporalMaxDeltaU$\%.
The requested spatial acceptance criterion is therefore \emph{not} uniformly satisfied.
\input{generated/table_bim_pfc_convergence.tex}

Table~\ref{tab:bimpfctarget} reports the distinct PFH cap-0.01 follow-up and the targeted
$N=128$ continuation of the out-of-screen KM $\alpha_{\mathrm c}=1$, $\gamma=4$ estimate.
Changing both $N$ and the cap from the original baseline is not a passing unchanged-step
spatial pair; the comparison between $N=64,96$ at cap 0.01 is instead a separate spatial test.
All four such PFH pairs complete with speed changes at most $\bimPFCTargetStableMaxDeltaU$\%.
For KM $\alpha_{\mathrm c}=1$, $\gamma=4$, the $N=96$ large-cap value
$\PFHkmOneNinetySixLargeStepU$\,m/s contrasts with $\PFHkmOneNinetySixSmallStepU$\,m/s
at cap 0.01, a change of $\PFHkmOneNinetySixCapDeltaU$\%. The sampled profile energy changes
$\PFHkmOneNinetySixCapDeltaEnergy$\% and its coverage fraction changes
$\PFHkmOneNinetySixCapDeltaCoverage$\% relatively. These expose sensitivity to the
integration control at that mesh.
The targeted $N=128$, cap-0.01 run aborts before impact. Its speed is unavailable and
excluded, so a finest $96\to128$ speed change cannot be computed and the finer spatial
convergence remains unconfirmed. No such refinement makes an out-of-screen incompressible
estimate physically admissible.
\input{generated/table_bim_pfc_target.tex}

Selected $N=96$ slow-condensation samples are $\PFPbimRefinedSlowSample$\,m/s (PFP) and
$\PFHbimRefinedSlowSample$\,m/s (PFH), compared with the original whole-grid largest samples
$\PFPenvSlowUmax$ and $\PFHenvSlowUmax$\,m/s. Speed stability does not imply profile stability:
the requested spatial changes reach $\bimPFCSpatialMaxDeltaMass$\% in sampled mass and
$\bimPFCSpatialMaxDeltaEnergy$\% in sampled kinetic energy. The exclusion of points near the
surface is resolution-dependent; coverage and separate profile changes are listed in
Table~\ref{tab:bimpfcprofiles} (Appendix~\ref{app:verif}). Thus the reported profile energies
and masses carry appreciably larger quadrature sensitivity than the far-pole speeds.

For the model-screen crossings, Table~\ref{tab:bimpfccross} gives the original interpolation,
best available $N=96$ estimate and empirical endpoint envelope. All use the broad sampled
stand-off bracket $[3,4]$; the much narrower numerical endpoint envelope is not a confidence
interval or a bound on interpolation error between samples. Mixed step settings after an
abort are identified explicitly. The physical fast-condensation maximum and continuous
benchmark optimum remain unresolved.
\input{generated/table_bim_pfc_crossings.tex}

\subsection{Step J6 --- Mass-resolved jet speed (observable O2)}
At impact the jet is the liquid enclosed by the re-entrant cavity surface between the tip and the rim. Its
axial velocity is evaluated on a $40\times10$ Gauss grid from the interior representation (points closer
than 0.6 element lengths to the surface are excluded; coverage varies by case and resolution), giving the
distribution $m_{\mathrm j}(\ge U)$ of Eq.~\eqref{eq:Um}. For the constant-pressure benchmark the jet
carries $\jetEfracMin$--$\jetEfracMax$\% of the cavity energy $\tfrac43\pi\Rm^3\Delta p$, nearly independent of
$\gamma$, consistent with the scale-free argument of Supponen et al.\ (mass $\propto\zeta^2$, speed
$\propto\zeta^{-1}$). The computed distribution predicts the mass-qualified speed
$U_{\mathrm j}(m_*)$ for its sampled profile. Independently, mass moving at speed at least $U$ carries
kinetic energy at least $m_*U^2/2$, so Eq.~\eqref{eq:Kj} implies the rigorous conditional
ceiling $U_{\mathrm j}(m_*)\le\sqrt{2\Phia/m_*}$ for $m_*>0$. Missing near-interface quadrature
mass and resolution changes are numerical uncertainties (Table~\ref{tab:bimpfcprofiles}); the distribution is not a certified
physical upper bound. Both comparisons appear in Fig.~\ref{fig:envelope}.

\subsection{Validity of the jet calculation}
(i) Potential flow neglects viscosity; at $L\approx24\,\mu$m the Reynolds number based on
$\sqrt{\Delta p/\rho}$ is $\approx130$, so viscous losses of order 10\% in the jet speed cannot be excluded
(not computed or bounded here). (ii) Incompressible: an adopted working screen is
$U_{\mathrm j}/c_{\mathrm w}\lesssim0.3$; it is not an error certificate even below the screen. Faster
predictions are flagged. (iii) The near-contact regime ($\gamma\lesssim0.2$ in Lechner et al.\ 2020) where
annular-collision needle jets of order 1000\,m/s form is excluded from $\mathcal D$; it requires
viscous compressible simulations and is not reachable by this method. (iv) The spherical-content closure
is an approximation (Section~\ref{sec:validity}). (v) Gel-embedded jets are not computed.
(vi) The reported far-pole maximum uses a $0.01L$ tip--wall stopping gap. It is not the physical
speed at contact; the potential-flow contact singularity cannot certify that endpoint.
